% ID: SEL-D0036-OA-E02
% ESTADO: cotejada. Segunda pasada visual y borrador compilado revisados.
% Fuente: Selectividad/matematicas_aplicadas/Ua1MATEMÁTICAS APLICADAS examen.pdf
% SHA-256: 77e0c295529ccd4ed502b5859d4011cb4b0e955a1e4795e570a9da20ae48293c
% Página: 1; opción A; ejercicio original 2.

Una entidad financiera lanza al mercado un plan de inversión cuya rentabilidad, $R(x)$, en miles de euros, viene dada por la función
\[
R(x)=-0.001x^2+0.5x+2.5\qquad1\leq x\leq500,
\]
donde $x$ es la cantidad de dinero invertida en miles de euros.
\begin{enumerate}[label=\alph*)]
\item \textbf{(1 punto)} Determine qué cantidad de dinero se debe invertir para obtener la máxima rentabilidad.
\item \textbf{(0.5 puntos)} \textquestiondown{}Qué rentabilidad se obtendría con dicha inversión?
\item \textbf{(1 punto)} \textquestiondown{}Cuál es la cantidad de dinero para la que se obtiene menor rentabilidad?
\end{enumerate}
